\begin{figure}[tbp]
\centering
\begin{tikzpicture}[x=0.86cm,y=0.88cm,>=Stealth,
 every node/.style={font=\small}]
 \definecolor{downpath}{RGB}{165,55,50}
 \definecolor{upfill}{RGB}{25,108,105}
 \begin{scope}
  \node[font=\small\bfseries] at (1.5,4.5) {Fixed isolated down strand};
  \foreach \x in {0,1,2,3}
   \draw[gray!40] (\x,0.25)--(\x,3.8);
  \draw[->,gray] (-0.65,0.35)--(-0.65,3.7) node[above] {$t$};
  \draw[downpath,very thick] (1,0.35)--(1,1.15)--(2,1.15)--(2,2.55)--(1,2.55)--(1,3.7);
  \draw[downpath,->,thick] (1.2,1.15)--(1.85,1.15);
  \draw[downpath,->,thick] (1.8,2.55)--(1.15,2.55);
  \fill[downpath] (1,0.35) circle (2pt);
  \fill[downpath] (1,3.7) circle (2pt);
  \node[left] at (0,0.35) {$t_0$};
  \node[left] at (0,3.7) {$t_1$};
  \node[right,align=left,text=downpath] at (2.2,1.6) {$X$};
  \node[align=center] at (1.5,-0.35) {The incident holding intervals\\contain no exchange marks.};
 \end{scope}
 \begin{scope}[xshift=6.8cm]
  \node[font=\small\bfseries] at (1.5,4.5) {Recolor, delete, and fill};
  \foreach \x in {0,1,2,3}
   \draw[upfill!45] (\x,0.25)--(\x,3.8);
  \draw[gray,dashed,thick] (1,0.35)--(1,1.15)--(2,1.15)--(2,2.55)--(1,2.55)--(1,3.7);
  \foreach \x/\y in {0/0.8,1/1.7,2/2.1,0/3.15}
   \draw[upfill,very thick] (\x,\y)--++(1,0);
  \draw[upfill,fill=white,thick] (1,0.35) circle (3pt);
  \draw[upfill,fill=white,thick] (1,3.7) circle (3pt);
  \node[right,fill=white,inner sep=2pt] at (1.1,0.35) {down $\to$ up};
  \node[right,fill=white,inner sep=2pt] at (1.1,3.7) {up $\to$ down};
  \node[align=center] at (1.5,-0.35) {New incident marks exchange\\two up colors.};
 \end{scope}
\end{tikzpicture}
\caption{The color-counting comparison for a fixed isolated strand.
Time runs upwards; horizontal segments represent instantaneous exchanges,
and the slot coordinates are schematic. Recoloring the prescribed segment
$X$ up allows its jump marks to be removed and its incident holding intervals
to be filled freely. The only remaining color discontinuities are the two
directed endpoint flips. Each cycle meeting a flip has at most one admissible
coloring, which gives the pin-weight bound used in the path comparison.}
\label{fig:surgery}
\end{figure}
